\chapter{Six problèmes de synthèse}
\section*{Problème 1 -- Un jardin et son coût}
Un jardin rectangulaire s'appuie sur un mur. On dispose de $40\unite m$ de grillage pour les trois autres côtés. On note $x$ la largeur perpendiculaire au mur.
\begin{enumerate}
\item Exprimer la longueur et donner le domaine géométrique de $x$.
\item Montrer que l'aire vaut $A(x)=200-2(x-10)^2$. Déterminer son maximum.
\item Déterminer les largeurs donnant une aire d'au moins $150\unite{m^2}$.
\item Deux entreprises proposent $120+15L$ et $200+12L$ euros pour poser $L$ mètres de grillage. Comparer les devis pour $L=40$, puis pour un $L\ge0$ quelconque.
\end{enumerate}

\section*{Problème 2 -- Un triangle sous plusieurs regards}
Dans un repère orthonormé, $A(0;0)$, $B(6;0)$ et $C(2;4)$. On note I le milieu de BC et G le point de coordonnées $(8/3;4/3)$.
\begin{enumerate}
\item Calculer les longueurs des côtés et l'aire.
\item Trouver une équation de la droite AI et montrer que G y appartient.
\item Vérifier que G appartient aussi à la médiane issue de B ; préciser sa position sur chaque médiane.
\item Trouver le centre O du cercle circonscrit et son rayon.
\item Déterminer le point d'intersection des hauteurs issues de C et de A.
\end{enumerate}
Les questions 3 à 5 prolongent facultativement la partie obligatoire.

\section*{Problème 3 -- Un tarif et une expérience}
Un appareil coûte $C(n)=80+0{,}25n$ euros pour $n$ utilisations, $n$ entier strictement positif. Son coût moyen est $m(n)=C(n)/n$.
\begin{enumerate}
\item Écrire $m(n)$ à l'aide de la fonction inverse et justifier son sens de variation.
\item Déterminer le plus petit $n$ tel que $m(n)\le1$ euro.
\item Un élève trace $m(x)$ sur $]0;200]$. Expliquer la différence entre ce tracé et les valeurs autorisées par le modèle.
\item Écrire une boucle cherchant le premier entier vérifiant le seuil, et prouver qu'elle se termine.
\end{enumerate}

\section*{Problème 4 -- Comparer deux distributions}
Deux groupes de $20$ élèves ont obtenu les notes suivantes : groupe A, dix $8$ et dix $12$ ; groupe B, cinq $4$, dix $10$ et cinq $16$.
\begin{enumerate}
\item Calculer moyenne, médiane, quartiles et écart-type dans chaque groupe.
\item Représenter les diagrammes en boîte et rédiger une comparaison.
\item Le groupe A gagne uniformément $2$ points ; le groupe B gagne $10\,\%$ de sa note. Calculer les nouvelles moyennes et dispersions.
\item Peut-on affirmer que les deux transformations profitent pareillement à tous les élèves ?
\end{enumerate}

\section*{Problème 5 -- Du contrôle qualité à la décision}
Une production comporte $3\,\%$ de pièces défectueuses. Un contrôle rejette $95\,\%$ des pièces défectueuses et $4\,\%$ des pièces conformes.
\begin{enumerate}
\item Construire un arbre puis un tableau pour $100\,000$ pièces théoriques.
\item Calculer la probabilité de rejet et la proportion de pièces défectueuses parmi les rejets.
\item Calculer la proportion de pièces défectueuses parmi les pièces acceptées.
\item Un second contrôle identique permet-il de multiplier automatiquement les deux probabilités de rejet ? Expliquer l'hypothèse manquante.
\end{enumerate}

\section*{Problème 6 -- Un atelier de preuve et de données}
On étudie les populations des douze cantons dans le fichier CSV fourni.
\begin{enumerate}
\item Vérifier les totaux femmes, hommes et ensemble.
\item Démontrer que la proportion nationale de femmes est une moyenne des proportions cantonales pondérée par les populations des cantons.
\item Expliquer pourquoi la moyenne simple des douze proportions ne donne pas forcément cette proportion nationale.
\item Programmer la construction du tableau Esch / autres cantons croisé avec femmes / hommes, puis vérifier ses marges.
\item \textbf{Facultatif.} Démontrer pour deux nombres $u,v$ et $0\le t\le1$ que $tu+(1-t)v$ est compris entre $u$ et $v$. Relier cela à la question 2.
\end{enumerate}
